\section{Outstanding analytical tasks}\label{sec:outstanding}

The proved chain ends at a uniform exterior enclosure and a collection
of localized and averaged moment estimates. Further progress requires
new quantitative input at identifiable points in that chain. The
following tasks are stated in the order of their immediate dependence.

\subsection{Acquire or test the every-center signed moment}

The nearest open positive task is to prove
\[
 Q(X)\le4/66^2\qquad\hbox{for every sufficiently large real }X,
\]
or a different complete boundary estimate below the Green margin in
\eqref{eq:margin-81}. Localization makes a uniformly strict cap on
$\|M_X\|_P$ sufficient after paying its norm error. The estimates in
Section~\ref{sec:means} allow rare large centers and therefore do not
complete this task. The missing quantity is a signed correlation with
the actual logarithmic phases, including nearby indices.

The complementary diagnostic task is the fixed short-window witness
$A_U(Y)\ge5/66^2$ along an unbounded sequence. Theorem~\ref{thm:quartic}
shows that one genuine alignment does not carry enough weighted mass.
Theorem~\ref{thm:sparse} further shows that any proposed mechanism
confined to $o(U^{1/64})$ of these logarithmic cells changes the full
moment by only $o(1)$. A new lower mechanism must control a more
distributed portion of the signed form, including its interaction
with the complement. It is not enough to display a constructive
subsum or to replace the complement by a growing upper allowance.
No such distributed mechanism is established in this manuscript.

\subsection{Extend the carrier and arithmetic control to deeper collars}

At the exact center on the curve,
\begin{equation}\label{eq:carrier-half}
 \Re s_{t_X}(X+\ii y)
 =\frac34+\frac y2+
          \frac{\log(1+(y/X)^2)}{8L_X}.
\end{equation}
The Euler-product margin is uniform for each fixed $y>1/2$.
At $y=1/2$, the remaining positive correction tends to zero;
below $1/2$ the real part is eventually below one. Thus the particular
elementary reciprocal-zeta proof used here loses its uniform margin
near $1/2$. This is a limitation of this carrier estimate, not a
theorem that every possible analytic carrier must fail.

An iteration needs a nonvanishing normalization or another signed
formulation on the actual deeper domains, together with complete
arithmetic error estimates there. The scale-invariant geometry of
Section~\ref{sec:geometry} does not supply either. Every stage must
also be available at the height under consideration: an unbounded
sequence of stage onsets would not justify infinitely many steps at
a fixed height.

\subsection{Pay the common-time compact region and the near-zero-time barrier}

The exterior theorem leaves $|x|<R(t)$, where $R(t)=4\pi\ee^{1/t}$.
At a fixed final time $T$, a direct zero exclusion on that full
compact complement would complete a global strip of the same width.
A different route is Polymath's three-part barrier criterion
\cite[Proposition 3.3]{Polymath}. To specify the missing work, fix
$0<y_0<1$, $0<T<1/2$ with $y_0^2+2T<1$, and set
$d_0=\sqrt{1-y_0^2}$. Its three zero-free regions are
\begin{align*}
 &t=0,\quad0\le x\le X_0,\quad
       \sqrt{y_0^2+2T}\le y\le1;\\
 &t=T,\quad x\ge X_0+d_0,\quad
       y_0\le y\le\sqrt{1-2T};\\
 &0\le t\le T,\quad X_0\le x\le X_0+d_0,\quad
       \sqrt{y_0^2+2(T-t)}\le y\le\sqrt{1-2t}.
\end{align*}
For $y_0=567/1000$, $T\le t_*$ and
$X_0+d_0\ge R(T)$, our theorem supplies the final exterior.
The initial region and the portion of the intermediate barrier with
$t<1/L_x$ remain outside its guarantee. In particular this includes
$t=0$. A complete certificate by either route would give a global
strip at $T$, and strip contraction would then imply
$\Lambda\le T+y_0^2/2$. No such certificate is supplied here.

The available onset is symbolic. Its sufficient construction already
requires $L_*\ge\kappa^{-1}=200000$, before the other constants are
evaluated. This describes the conservative proof recipe, not the true
first height at which the enclosure holds. A useful finite-height
bridge would require substantially sharper effective constants and
overlap with separately established compact control. The earlier
global-use audit's numerical comparison at width $4/7$ is specific
to that width~\cite{BB16}; it is not automatically a comparison for
$567/1000$.

\subsection{Justify any time-zero passage or iteration}

For every fixed finite $x$, the condition $|x|\ge4\pi\ee^{1/t}$
fails when $t$ is sufficiently close to zero. Hence the exterior
theorem cannot simply be sent to $t\downarrow0$ at that height.
Similarly, the conditional $9/16$ result is a curve statement;
it has no proved common-time extension. A limiting argument must
provide control on the relevant fixed compact sets, or a separate
uniform time-space propagation theorem with all boundary payments.

Finally, the complete corrected initial expansion and the joined
analytical chain merit specialist review independent of the shared
LLM source lineage. This would validate the foundations and the
new synthesis at a different evidence level. It would not by itself
fill the open moment, carrier or compact-region estimates. Those
remain the concrete mathematical tasks needed to turn the present
exterior progress into a global or iterative result.
